% !TeX root = ../exam-zh-doc-basic.tex

\section{完整示例}

本章提供可直接使用的完整试卷示例。所有示例代码都可以在 \file{examples-basic/} 目录中找到。

\subsection{示例文件说明}

\begin{description}
  \item[\file{00-minimal.tex}] 最小示例（10行代码创建试卷）
  \item[\file{01-first-exam.tex}] 第一份完整试卷（包含选择题、填空题、解答题）
  \item[\file{02-math-basic.tex}] 高中数学期末考试（带草稿纸）
  \item[\file{03-real-exam-scenario.tex}] 真实考试场景（带密封线、完整答案、评分标准）
\end{description}

\begin{tcolorbox}[colback=blue!5!white, colframe=blue!75!black, title=\textbf{如何使用这些示例}]
\textbf{方法一：直接编译}
\begin{enumerate}
  \item 进入 \file{examples-basic/} 目录
  \item 选择一个示例文件，如 \file{01-first-exam.tex}
  \item 使用 XeLaTeX 编译：\texttt{xelatex 01-first-exam.tex}
  \item 查看生成的 PDF 文件
\end{enumerate}

\textbf{方法二：复制修改}
\begin{enumerate}
  \item 复制示例文件到你的工作目录
  \item 修改标题、科目、题目内容
  \item 注释或取消注释 |\examsetup| 切换学生版/教师版
  \item 编译生成你自己的试卷
\end{enumerate}

\textbf{推荐：}从 \file{01-first-exam.tex} 开始，这是最适合新手的模板。
\end{tcolorbox}

\subsection{示例对比}

\begin{center}
\begin{tabular}{lp{5cm}p{5cm}}
\hline
\textbf{示例} & \textbf{特点} & \textbf{适用场景} \\
\hline
00-minimal & 最简代码，无答案控制 & 快速测试、学习基础语法 \\
01-first-exam & 完整结构，有答案控制 & 日常测验、课堂小考 \\
02-math-basic & 带草稿纸，题目较多 & 期中/期末考试 \\
03-real-exam & 带密封线、评分标准 & 正式考试、月考 \\
\hline
\end{tabular}
\end{center}

\subsection{数学试卷示例}

完整代码见 \file{examples-basic/01-first-exam.tex}，这里展示主要结构：

\begin{latexcode}[deletetexcs={\documentclass},morekeywords={\documentclass}]
  \documentclass{exam-zh}

  % 教师版配置（注释掉生成学生版）
  % \examsetup{
  %   question/show-answer = true,
  %   solution/show-solution = show-stay
  % }

  \begin{document}

  % 试卷抬头
  \title{2024年春季期中考试}
  \subject{数学}
  \maketitle

  \information{
    姓名\underline{\hspace{6em}},
    班级\underline{\hspace{6em}},
    座号\underline{\hspace{4em}}
  }

  \begin{notice}
    \item 本试卷共3大题，满分100分，考试时间90分钟。
    \item 请用黑色签字笔在答题卡上作答。
  \end{notice}

  % 一、选择题
  \section{选择题（每题5分，共20分）}
  \begin{question}
    已知集合 $A = \{1, 2, 3\}$，$B = \{2, 3, 4\}$，
    则 $A \cap B = $ \paren[B]
    \begin{choices}
      \item $\{1\}$
      \item $\{2, 3\}$   % 正确答案
      \item $\{1, 2, 3, 4\}$
      \item $\varnothing$
    \end{choices}
  \end{question}

  % 二、填空题
  \section{填空题（每题5分，共20分）}
  \examsetup{question/index = 0}  % 重置题号

  \begin{question}
    若 $f(x) = 3x - 2$，则 $f(2) = $ \fillin[4]。
  \end{question}

  % 三、解答题
  \section{解答题（共60分）}
  \examsetup{question/index = 0}

  \begin{problem}[points = 15]
    已知函数 $f(x) = x^2 - 4x + 3$。
    \begin{enumerate}
      \item 求 $f(x)$ 的对称轴；
      \item 求 $f(x)$ 的最小值。
    \end{enumerate}

    \begin{solution}
      \begin{enumerate}
        \item 对称轴为 $x = 2$。\score{7}
        \item 最小值为 $-1$。\score{8}
      \end{enumerate}
    \end{solution}
  \end{problem}

  \end{document}
\end{latexcode}

\subsection{真实考试场景示例}

完整代码见 \file{examples-basic/03-real-exam-scenario.tex}，展示了真实考试的完整要素：

\textbf{特色功能：}
\begin{enumerate}
  \item \textbf{密封线：}使用 |\secret| 命令生成左侧密封线
  \item \textbf{详细注意事项：}符合正式考试要求的考生须知
  \item \textbf{完整评分标准：}每个给分点都用 |\score{分值}| 标记
  \item \textbf{草稿纸：}使用 |\draftpaper| 命令生成方格纸
\end{enumerate}

\textbf{代码片段：}

\begin{latexcode}
  % 密封线
  \secret

  % 试卷信息
  \information{
    姓名\underline{\hspace{6em}},
    班级\underline{\hspace{6em}},
    学号\underline{\hspace{6em}},
    考场\underline{\hspace{4em}},
    座号\underline{\hspace{3em}}
  }

  % 详细注意事项
  \begin{notice}
    \item 本试卷共 4 页，22 小题，满分 150 分。
    \item 答卷前，考生务必将自己的姓名、学号填写在答题卡上。
    \item 作答选择题时，选出答案后，用2B铅笔在答题卡上涂黑。
    \item 非选择题必须用黑色字迹的钢笔或签字笔作答。
  \end{notice}

  % ... 题目内容 ...

  % 草稿纸（在新页面）
  \newpage
  \section*{草稿纸}
  \draftpaper
\end{latexcode}

\subsection{自定义你的试卷}

基于这些示例，你可以轻松自定义：

\begin{enumerate}
  \item \textbf{修改标题和科目}
    \begin{latexcode}
      \title{你的考试名称}
      \subject{你的科目}
    \end{latexcode}

  \item \textbf{调整题型分值}
    \begin{latexcode}
      \section{选择题（每题5分，共40分）}
      % 或
      \section{选择题（本题共8小题，每小题5分，共40分）}
    \end{latexcode}

  \item \textbf{添加或删除题目}

    直接复制 |question|、|problem| 环境块，修改内容即可

  \item \textbf{切换学生版/教师版}

    注释或取消注释文档开头的 |\examsetup| 配置
\end{enumerate}

\subsection{更多示例}

\textbf{官方仓库：}
\begin{itemize}
  \item GitHub: \url{https://github.com/xkwxdyy/exam-zh}
  \item 查看 \file{examples-basic/} 目录获取最新示例
  \item 查看 \file{example-single.tex} 获取高考真题示例
\end{itemize}

\textbf{社区贡献：}
\begin{itemize}
  \item 加入 QQ 交流群：652500180
  \item 分享你的试卷模板
  \item 获取更多使用技巧
\end{itemize}
  \begin{question}
    ...
  \end{question}

  % 三、解答题（4道）
  \section{解答题（共40分）}
  \examsetup{question/index = 0}
  \begin{problem}[points = 10]
    ...
    \begin{solution}
      ...
    \end{solution}
  \end{problem}

  % 草稿纸
  \newpage
  \section*{草稿纸}
  \draftpaper

  \end{document}
\end{latexcode}

\subsection{使用建议}

\subsubsection{如何使用示例}

\begin{enumerate}
  \item \textbf{复制示例文件}：
    \begin{itemize}
      \item 找到最接近你需求的示例
      \item 复制整个文件
      \item 重命名为你的试卷名称
    \end{itemize}

  \item \textbf{修改内容}：
    \begin{itemize}
      \item 修改标题和科目
      \item 替换题目内容
      \item 调整题目数量和分值
    \end{itemize}

  \item \textbf{生成PDF}：
    \begin{itemize}
      \item 使用 XeLaTeX 编译
      \item 检查效果
      \item 根据需要调整
    \end{itemize}
\end{enumerate}

\subsubsection{常见修改}

\textbf{修改题目数量：}

直接增加或删除 |question|/|problem| 环境即可。

\textbf{修改分值分配：}

修改每题的 |[points = ...]| 参数，并调整标题中的总分。

\textbf{切换学生版/教师版：}

注释或取消注释文档开头的 |\examsetup| 配置。

\textbf{调整页面布局：}

参考第4章"常见场景"的页面设置部分。

\subsection{代码模板}

这里提供一些常用的代码片段，可以直接复制使用。

\subsubsection{选择题模板}

\begin{latexcode}
  \begin{question}
    题干内容 \paren[A]
    \begin{choices}
      \item 选项A（正确答案）
      \item 选项B
      \item 选项C
      \item 选项D
    \end{choices}
  \end{question}
\end{latexcode}

\subsubsection{填空题模板}

\begin{latexcode}
  \begin{question}
    题干内容，答案是 \fillin[答案] 。
  \end{question}
\end{latexcode}

\subsubsection{解答题模板}

\begin{latexcode}
  \begin{problem}[points = 10]
    题干内容

    \begin{solution}
      解：

      步骤1 \score{3}

      步骤2 \score{4}

      答案 \score{3}
    \end{solution}
  \end{problem}
\end{latexcode}

\subsection{进一步学习}

\begin{itemize}
  \item 尝试修改示例文件，理解各部分的作用
  \item 参考第3章学习更多题型和功能
  \item 遇到问题查阅第6章FAQ
  \item 需要高级功能时查阅完整文档 \file{exam-zh-doc.pdf}
\end{itemize}
